\chapter{Needs and Strategies}

Feelings are rapid signals that tell us whether our needs are being met;
they serve as the bridge to understanding our inner world.
However, a feeling is merely a signal.
What truly drives this signal is the \textbf{need} behind it.
If feelings are the color and brightness of a light, then needs are the
energy source generating that light.
This chapter explores the essence of needs, the variety of universal human
needs, and how we can authentically identify and express our deepest needs
so that our communication genuinely touches the heart.

\section{Needs}

\subsection{Definition}

\subsubsection{The Source of Feelings}

Needs are the internal driving forces that we all share, which generate our
feelings.
Stepping away from dualistic thinking is also known as ``needs-based thinking.''
Feelings are not directly caused by the actions of others, but by whether
our own needs are being satisfied.
This is one of the core insights of Nonviolent Communication:
\textbf{No one can make you feel a certain way; your feelings come only
from your own needs} \citep{rosenberg2003}.
We take responsibility for our own needs, rather than pushing the
responsibility onto others or forcing ourselves to meet others' needs.
Let us return to the late-night homecoming scenario to
deeply understand the relationship between feelings and needs:

\paragraph{Case Analysis: The Husband Returning Late (Needs Version)}

The husband opens the door and comes home at 11 PM.
The wife has been waiting for a long time, and complex emotions are surging
within her.
She feels lonely and overwhelmed.
Is the husband's late return the ``cause'' of her loneliness?

On the surface, it seems so.
But if we ask a deeper question: \textbf{``She is lonely, because what need
is not being met?''}
The answer becomes clear: She is lonely because she deeply \textbf{needs
companionship, connection, and support} \citep{bowlby1969}.

This distinction is crucial:
\begin{itemize}
    \item \textbf{Judgmental Thinking (Alienating Communication):}
    ``You make me feel lonely, it's your fault, it's all because you didn't
    fulfill your duties as a husband.''
    (Pushing the responsibility for feelings onto the other person)
    When the exhausted husband pushes open the door, if he is met not with
    honesty about feelings, but with a cold accusation: ``You finally know
    how to come back? You don't have this family in your heart at all,
    you're too selfish!''

    These words act like an invisible subpoena, instantly dragging the
    husband to the \textbf{``defendant's seat''} in the relationship court.
    At this moment, the amygdala in the husband's brain sounds the alarm,
    activating a strong \textbf{defense mechanism} \citep{lieberman2007}.
    Sitting in the ``defendant's seat,'' he no longer has the capacity to
    notice his wife's loneliness, because he is busy searching for evidence
    to defend himself:
    ``Who am I working so late for? If I don't work hard overtime, who will
    pay the mortgage? You only see that I come back late, why don't you see
    how much I've given for this family?''

    In this trial of ``who is right and who is wrong,'' both sides try to
    prove that they are the ``victim'' and the other is the ``perpetrator.''
    As the debate intensifies, the wife continues: ``You still dare to
    argue? You don't have the ability to do your job well and you turn
    around and blame me? You don't deserve to be the man of this house!
    Chores, children, I'm doing it all! What have you done?''
    The husband is pushed to his limit and finally explodes: ``You don't
    deserve to be the woman of this house! What kind of family raised a
    shrew like you?''

    If the husband attacks outward, he might physically strike his wife.
    At this moment, conjugal affection turns into mutual hatred.
    The wife truly pushes the husband into the ``defendant's seat,'' and
    the husband goes to prison for domestic violence.
    The wife then ``no longer believes in love,'' raising the child alone.
    And the child, witnessing all this, grows up full of fear and insecurity,
    amidst the mother's curses and resentment toward the husband \citep{vanderkolk2014}.

    If the husband chooses to attack inward, he might choose jackal silence,
    or silently blame himself: ``Why am I so useless? I can't even support a
    family, and I get scolded like this by my wife? I really am a failure!''
    Physically and mentally exhausted, frustrated inside and out, he chooses
    to escape this relationship, leaving this home forever, and eventually
    sinks into illness and bad habits.
    The wife also feels deeper loneliness and despair due to his departure,
    developing psychological illnesses.
    The child, unable to go to school due to lack of money, blames the
    family's incompetence and heartlessness, eventually turning to crime \citep{levine1997}.

    \item \textbf{Empathetic Thinking (Natural Communication):}
    ``I feel lonely, because I really need your companionship.''
    (Stating the need behind the feeling)

    The husband pushes open the door, sees the tears shimmering in his wife's
    eyes, and hears her trembling voice say: ``I feel lonely, because I really
    need your companionship.''
    This sentence is like a key, unlocking each other's hearts.
    The husband holds his wife, and they weep in each other's arms. In this
    embrace, the wife feels a long-lost sense of security and being seen, and
    the husband also feels that he is satisfying his wife's needs bit by bit \citep{porges2011}.

    After calming down, the husband confesses to his wife the stress and
    exhaustion of overtime: ``Recently the company is laying off people, I'm
    under a lot of pressure, afraid of losing our source of income after
    losing my job, so I've been working late. I hope our child can realize
    their dreams, and we can also provide a high-quality life for our family.''
    Hearing her husband's honesty, the wife feels his vulnerability and effort,
    her heart filled with emotion and understanding.
    The two begin to discuss how to jointly meet the need for family financial
    security while also meeting the wife's need for companionship.
    Both sides agree that the husband will come home punctually at 7 PM on
    Mondays and Fridays, and work overtime at other times.

    The husband's company later laid off employees, but because of his
    previous hard work and performance, he successfully transferred to another
    department, continuing to provide stable income for the family.
    Because of her husband's understanding and support, the wife also feels
    seen and valued in this family, her heart filled with security and
    satisfaction.
    The child successfully entered their ideal university and found their
    passion.
    After experiencing this crisis, the family relationship became more solid
    and deep, empathizing and caring for each other even in difficulties.
    Influenced by their parents, the child later formed another family full
    of love, continuing to pass on love and happiness \citep{siegel2003}.
\end{itemize}

Happy families have not avoided misfortune, but when misfortune occurs,
both sides can return to an awareness of their own feelings and needs,
express them honestly, and find solutions that willingly satisfy each other
in the other's needs.
And our way of thinking in the face of misfortune determines different life
trajectories: hate and love are two expressions of the same encounter, and
which mindset we choose determines whether we move toward hate and violence
or toward love and peace \citep{gottman1997}.

\paragraph{The Commonality of Humanity}
Needs are the internal driving forces everyone possesses.
Regardless of race, gender, age, or background, everyone has the same core
needs.
We can imagine an extreme contrast: a CEO in a bustling metropolis and a
shepherd in a remote mountain area.
Their life paths may never cross, and their ways of satisfying needs are
worlds apart---the CEO may gain ``a sense of achievement'' by improving
performance, while the shepherd may gain ``a sense of security'' by
guarding the flock.
However, when we strip away these specific behaviors and touch the deepest
core, you will find that their needs for life are exactly the same; what
differs are the methods of realization \citep{ekman1992}.

This ``commonality'' is the foundation that makes empathy possible.
The reason we can empathize with others' pain is that we have the same
needs in our hearts.
When we can penetrate the other person's angry expressions, accusing words,
and seemingly absurd behavior, and see the universal human needs behind
them that are no different from our own, the other person is no longer an
``unreasonable freak'' or ``enemy,'' but a person just like us who is
trying hard to find satisfaction.
Seeing commonality is the most effective antidote to dissolving ``hostility.''
Focusing on universal human needs is the ``wick'' for us to detach from
dualistic thinking.

\begin{keyinsight}
Our feelings are not caused by the actions of others, but by our own
underlying needs being met or unmet.
By recognizing that all human beings share the same core needs, we can
transform conflict into connection and foster empathy, completely shifting
from judgmental defense to compassionate understanding.
\end{keyinsight}
